%%%%%%%%%%  scan idx 197  |  folio 184  %%%%%%%%%%
\textbf{Proposition 2.4.8.}\footnote{Ce résultat, ainsi que sa démonstration, a été communiqué gracieusement au rédacteur par A. Grothendieck ; il a été obtenu indépendamment par J. Frish [5].}. Soit $S$ un espace pseudo-analytique complexe, compact, et de Stein (2.4.7). On suppose que, pour tout ouvert $U$ de $S$ et tout sous-espace fermé réduit $X$ de $U$, le normalisé de $X$ est localement connexe\footnote{Nous admettrons que les notions de Géométrie analytique locale de [10] s'étendent aux espaces pseudo-analytiques complexes ; la réduction au cas habituel est en effet chaque fois immédiate.} (condition vérifiée par exemple si $S$ est défini localement par des équations ou inéquations analytiques réelles dans un espace $\mathbb R^n$ [15]). Alors l'anneau $A=\Gamma(S,\mathcal O_S)$ est un anneau noethérien.

\emph{Preuve.} Il s'agit de montrer que, si $M$ est un $A$-module de présentation finie, l'ensemble des sous-modules (ou, ce qui revient au même, l'ensemble des sous-modules de type fini) de $M$ vérifie la condition maximale. Or, si $N$ est un sous-module de type fini de $M$, alors $M/N$ est de présentation finie d'après (I 2.6 et 2.9), donc $N$ est de présentation finie d'après l'équivalence (2.4.3 c)). Utilisant à nouveau l'équivalence (2.4.3 c)) entre la catégorie des $\mathcal O_S$-Modules cohérents et la catégorie des $A$-modules de présentation finie, et la compacité de $S$, on est ramené à prouver le lemme suivant :

\textbf{Lemme 2.4.8.1.} Soit $S$ un espace pseudo-analytique complexe vérifiant la condition de locale connexité de (2.4.8). Soient $F$ un $\mathcal O_S$-Module cohérent, $(E_n)$ une suite croissante de sous-$\mathcal O_S$-Modules cohérents de $F$, $s$ un point de $S$. Il existe alors un voisinage ouvert $U$ de $s$ tel que la suite $(E_n|U)$ soit stationnaire.

\emph{Preuve.} L'anneau $\mathcal O_{S,s}$ étant noethérien, il existe un entier $n_0$ tel $(E_n)_s=(E_{n_0})_s$ pour $n\geq n_0$. Quitte à remplacer la suite $(E_n)$ par la suite $(E_n/E_{n_0})$ pour $n\geq n_0$, on peut donc supposer que $(E_n)_s=0$ pour tout $n$. D'autre part, $F_s=M$ est un $\mathcal O_s$-module de type fini, donc il existe (Bourbaki Alg. com. chap. IV \S1 th.1) une filtration finie
\[
 M=M_0\supset M_1\supset\ldots\supset M_q=0
\]
telle que $M_i/M_{i+1}\simeq \mathcal O_s/p_i$, où $p_i$ est un idéal premier de $\mathcal O_s$. Comme la catégorie des $\mathcal O_s$-modules de type fini est équivalente à la catégorie des germes en $s$ de $\mathcal O_S$-Modules cohérents, on peut supposer, quitte à se
loca%
